\scanpage{61}
% archive sheet number: 60
$\unsure{I|\mathcal I}$ ensemble ordonné satisfaisant à ``condition des chaines''
\margin{a) La longueur des chaines \unsure{joignant} $x\geq y$ est bornée\par b) Deux chaines saturées joignant $x$ à $y$ ont $=$ longueur soit $d(x,y)$}

Pour toute chaine \uline{saturée} $c=(x_0\geq x_1\geq\cdots\geq x_n)$, supposons donné un \uline{anneau}
\[\begin{aligned}\longrightarrow\\A_c=A_{[x_0,\ldots,x_n]}\end{aligned}\]
et pour une chaine de longueur \unsure{non nulle|supérieure} (supposons $n\geq1$) des \uline{hom d'anneaux}
\[
\begin{aligned}
u_1:\quad&A_{[x_0,\ldots,x_{n-1}]}\longrightarrow A_{[x_0,\ldots,x_n]}\\
u_2:\quad&A_{[x_1,\ldots,x_n]}\longrightarrow A_{[x_0,\ldots,x_n]}
\end{aligned}
\]
\margin{[c'est $(1_{x_n}\Theta)\circ$]\par [c'est $\circ(\Theta1_{x_0})$]}

[de sorte que $u_1,u_2:A_{c'}\longrightarrow\displaystyle\prod_{\substack{c\supseteq c'\\\mathrm{lg}\,c=\mathrm{lg}\,c'+1}}A_c$]
\[\deleted{\unsure{u_2:A_{c'}\longrightarrow\prod_{c\not\supset c'} A_c}}\]

on suppose que $u_1$ et $u_2$ commutent: si $n\geq2$
\diagram{\begin{tikzcd}[ampersand replacement=\&]
A_{[\unsure{x_0|x_1},\ldots,x_{n-1}]}\arrow[r,"u_2"]\arrow[d,"u_1"']\&A_{[x_0,x_1,\ldots,x_{\unsure{n-1|n}}]}\arrow[d,"u_1"]\\
A_{[x_1,\ldots,x_{\unsure{n-1|n}}]}\arrow[r]\&A_{[x_0,x_1,\ldots,x_{n-1},x_n]}
\end{tikzcd}}

de sorte que cela permet de définir sans ambiguïté des hom d'anneaux
\[u_{c,c'}\quad A_c\longrightarrow A_{c'}\qquad\text{si }c\subset c'\]
\deleted{[\unsure{donc}} et on a une condition de \uline{transitivité}, qui exprime que $(A_c)$ est un système inductif sur l'ens des chaines ordonnées. De ceci, et de la loi \uline{\unsure{multiplicative|multiplication}} dans les $A_c$, on déduit des \deleted{\unsure{hom|homomorphismes}} \inserted{accouplements \unsure{bi-additifs|biadditifs}}
\[A_{c'}\times A_c\longrightarrow A_{c'c}\]
si $c',c$ sont composables dans cet ordre, et cette loi
